\documentclass[../../../include/open-logic-section]{subfiles}

\begin{document}

\olfileid[tr]{sth}{card-arithmetic}{fix}
\olsection{$\aleph$-Sabit Noktaları}

\olref[spine][]{chap} içinde Yerine Koyma'nın gerekçelendirilmesi
gerektiğini; çünkü hiyerarşiyi epey yüksek olmaya zorladığını öne sürdük.
Artık biraz kardinal aritmetiği yaptığımıza göre hiyerarşinin yüksekliğine
ilişkin küçük bir örnek verebiliriz.

Açıkça $0<\aleph_0$, $1<\aleph_1$, $2<\aleph_2$\ldots olur ve gerçekten
büyüklük farkı her adımda yalnızca \emph{artar}. Bu yüzden her $\kappa$
sıra sayısı için $\kappa<\aleph_\kappa$ olduğunu öne sürmek çekicidir.

Fakat $\ZFC$ kabul edildiğinde bu varsayım \emph{yanlıştır}. Gerçekten
$\kappa=\aleph_\kappa$ olan $\kappa$ kardinal sayılarının, yani
\emph{$\aleph$-sabit noktalarının} bulunduğunu ispatlayabiliriz.

\begin{prop}\ollabel{alephfixed}
Bir $\aleph$-sabit noktası vardır.
\end{prop}

\begin{proof}
Özyineleme kullanarak şunları tanımlayın:
\begin{align*}
	\kappa_0 &= 0\\
	\kappa_{n+1} &= \aleph_{\kappa_n}\\
	\kappa&= \bigcup_{n < \omega}\kappa_n
\end{align*}
Şimdi $\kappa$, \olref[cardinals][classing]{unioncardinalscardinal}
gereği bir kardinal sayıdır. Üstelik:
\[
	\kappa= \bigcup_{n < \omega} \kappa_{n+1} =
	\bigcup_{n < \omega}\aleph_{\kappa_n} =
	\bigcup_{\alpha < \kappa}\aleph_\alpha = \aleph_\kappa
\]
% By construction, $\kappa$ is the least cardinal greater than each
% $\kappa_n$. So $\aleph_\kappa$ is the least cardinal greater than
% each $\aleph_{\kappa_n}$, and hence greater than each
% $\kappa_{n+1}$. But equally $\kappa$ is the least cardinal greater
% than each $\kappa_{n+1} = \aleph_{\kappa_n}$. So $\kappa=
% \aleph_\kappa$.
\end{proof}

Boolos bir zamanlar tam da az önce kurduğumuz $\aleph$-sabit noktası
hakkında bir makale yazdı. Makalesinin başında $\kappa$'nın varlığını
belirttikten sonra şöyle dedi:
\begin{quote}
	[$\kappa$], küme teorisiyle daha önce karşılaşmamış kişilerin ölçülerine
	göre \emph{epey büyük} bir sayıdır; bana öyle görünüyor ki, o kadar
	büyüktür ki iddialarından biri en az $\kappa$ tane nesne bulunduğu olan
	her teorinin doğruluğunu sorgulamayı gerektirir.
	\citep[p.~257]{Boolos2000}
\end{quote}
Makalesini sonunda şu soruyla bitirdi:
\begin{quote}
	[öykünün] başında durum her ne olmuş olursa olsun, buraya kadar
	geldiğimizde çarkların boşuna döndüğünden ve artık gerçekte olan herhangi
	bir şeyin betimlemesini dinlemediğimizden kuşkulanıyor muyuz?
	\citep[p.~268]{Boolos2000}
\end{quote}
Gerçekten ``gerçekte olan herhangi bir şeyi'' aşmışsak suçlayacağımız yer
doğrudan Yerine Koyma olmalıdır. Çünkü ispatımızın işlemesine izin veren bu
aksiyomdur. Öyleyse Boolos'un birkaç on yıl önce Yerine Koyma'nın ``hiçbir
istenmeyen'' sonucu olmadığı yönündeki iddiasını yeniden değerlendirmesi
gerekeceği varsayılır (bkz. \olref[replacement][extrinsic]{sec}).

Fakat $\kappa$'nın varlığı o kadar kötü müdür? Burada Russell'ın
\emph{Tristram Shandy paradoksunu} ele almak yararlı olabilir. Tristram
Shandy yaşamını günlüğünde belgeler, fakat tek bir günü kaydetmesi bir yıl
sürer. Tristram geçen her yılla daha da geride kalır: bir yıl sonra yalnızca
bir günü kaydetmiş, 364 günü kaydetmeden yaşamıştır; iki yıl sonra yalnızca
iki günü kaydetmiş, 728 günü kaydetmeden yaşamıştır; üç yıl sonra yalnızca
üç günü kaydetmiş, 1092 günü kaydetmeden yaşamıştır
\dots\footnote{Artık yılları göz ardı ediyoruz.} Yine de Tristram
\emph{ölümsüzse} her günü kaydetmeyi başarır; çünkü yaşamının $n$inci
yılında $n$inci günü kaydedecektir. Dolayısıyla ``zamanın sonunda''
Tristram'ın eksiksiz bir günlüğü olacaktır.

Öyleyse şu düşüncenin bundan farkı nedir: $\alpha$, $\aleph_\alpha$'dan
küçüktür---hatta gittikçe umutsuzlaşan biçimde daha küçüktür---ta ki
$\kappa$'ya kadar; bu noktada arayı kapatırız ve
$\kappa=\aleph_\kappa$ olur?

Bunu bir yana bırakıp $\ZFC$'yi kabul ettiğimizi varsayarak sabit nokta
kuruluşlarına ilişkin biraz daha eğlenceyle bitirelim. Sıradaki üç sonuç,
sezgisel olarak hiyerarşinin yüksek olduğu kadar geniş olduğu (önemsiz
olmayan) bir nokta bulunduğunu kurar:

\begin{prop}\ollabel{bethfixed}
Bir $\beth$-sabit noktası, yani $\kappa=\beth_\kappa$ olan bir $\kappa$
vardır.
\end{prop}

\begin{proof}
\olref{alephfixed} içindeki gibi, ``$\aleph$'' yerine ``$\beth$''
kullanın.
\end{proof}

\begin{prop}\ollabel{stagesize}
$\card{V_{\omega+\alpha}}=\beth_\alpha$. Eğer
$\omega\ordtimes\omega\leq\alpha$ ise
$\card{V_\alpha}=\beth_\alpha$.
\end{prop}

\begin{proof}
İlk iddia basit bir sonlu ötesi tümevarımla çıkar. İkinci iddia ise
$\omega\ordtimes\omega\leq\alpha$ olduğunda
$\omega+\alpha=\alpha$ olmasından çıkar. Bunu kurmak için
\olref[ord-arithmetic][]{chap} içindeki sıra sayısı aritmetiği olgularını
kullanırız. Önce
$\omega\ordtimes\omega=\omega\ordtimes(1\ordplus\omega)=
(\omega\ordtimes1)\ordplus(\omega\ordtimes\omega)=
\omega\ordplus(\omega\ordtimes\omega)$ olduğuna dikkat edin. Şimdi
$\omega\ordtimes\omega\leq\alpha$, yani bir $\beta$ için
$\alpha=(\omega\ordtimes\omega)\ordplus\beta$ ise
$\omega\ordplus\alpha=\omega\ordplus
((\omega\ordtimes\omega)\ordplus\beta)=
(\omega\ordplus(\omega\ordtimes\omega))\ordplus\beta=
(\omega\ordtimes\omega)\ordplus\beta=\alpha$ olur.
\end{proof}

\begin{cor}
$\card{V_\kappa}=\kappa$ olan bir $\kappa$ vardır.
\end{cor}

\begin{proof}
$\kappa$, \olref{bethfixed} tarafından verilen bir $\beth$-sabit noktası
olsun. Açıkça $\omega\ordtimes\omega<\kappa$. Dolayısıyla
\olref{stagesize} gereği
$\card{V_\kappa}=\beth_\kappa=\kappa$ olur.
\end{proof}

$V_\kappa$'nın altında kaç aşama varsa $V_\kappa$'nın o kadar
!!{element}s vardır. Öyleyse sezgisel olarak $V_\kappa$ yüksek olduğu kadar
geniştir. Bu, Tristram Shandy'yi çok andırır: \emph{kuvvet kümeleri} alarak bir
aşamadan ötekine geçer, böylece hiyerarşimizi her adımda \emph{çok} daha
büyük yaparız. Fakat ``sonunda'', yani $\kappa$ aşamasında, hiyerarşinin
genişliği yüksekliğine yetişir.

Şu sorulabilir: \emph{Hiyerarşinin genişliği yüksekliğiyle ne sıklıkta
eşleşir?} Yanıt şudur: \emph{Sıra sayıları ne kadar sıksa o kadar sık.}
Fakat bunun biraz açıklanması gerekir.

Bir $\tau$ terimini şöyle tanımlarız. Her $A$ için:
\begin{align*}
	\tau_0(A) & \defis \card{A}\\
	\tau_{n+1}(A) & \defis \beth_{\tau_n(A)}\\
	\tau(A) & \defis \bigcup_{n < \omega}\tau_n(A)
\intertext{Her $A$ için $\tau(A)$, \olref{bethfixed} içindeki gibi bir
$\beth$-sabit noktasıdır ve açıkça $\card{A}<\tau(A)$ olur. Şimdi şu
özyinelemeli tanımı ele alın:}
	W_0 &\defis 0\\
	W_{\alpha + 1} & \defis \tau(W_\alpha)\\
	W_\alpha & \defis \bigcup_{\beta < \alpha} W_\beta
	\text{, $\alpha$ bir limit olduğunda}
\end{align*}
Kuruluş bütün sıra sayıları için tanımlıdır. Öyleyse sezgisel olarak $W$,
sıra sayılarından $\beth$-sabit noktalarına giden bir
``!!a{injection}''dur. Ve tam önceki gibi her $\alpha$ için
$V_{W_\alpha}$ yüksek olduğu kadar geniştir.

\end{document}
